% =====================================================================
% preamble/macros.tex
% 目的：
%  - 図表・式などの参照コマンドを統一
% 前提：
%  - styles/thesis-packages.sty 側で graphicx / subfiles を読み込み済み
%  - 本プロジェクトは dvipdfmx（.pdf/.png/.jpg の画像を推奨）
% =====================================================================

% ---------------- 参照系（Fig., Table, Eq. 等を統一） --------------------
% 使い方： \figref{fig:setup}, \tabref{tab:params}, \eqnref{eq:newton}, \chapref{chap:results}, \secref{sec:methods}, \appref{app:hardware}
% 出力例： "Fig. 1", "Table 2", "Eq. (3)", "3章", "2.3節", "付録A"
\newcommand{\figref}[1]{Fig.~\ref{#1}}
\newcommand{\tabref}[1]{Table~\ref{#1}}
\newcommand{\eqnref}[1]{式(\ref{#1})}
\newcommand{\chapref}[1]{\ref{#1}章}
\newcommand{\secref}[1]{\ref{#1}節}
\newcommand{\appref}[1]{付録\ref{#1}}

% ---------------- 画像パス設定（subfiles対応） -------------------------
% \graphicspath + \subfix により、親(main.tex)でも子(chapX.tex)でも
% 「figures/」配下を自動で探索します。
% 例：\includegraphics{chap1/sample_fig.pdf} だけでOK。
% 重要：
%  - 図は figures/chapX/ に置く運用を推奨
\graphicspath{{\subfix{figures/}}}

% ---------------- よく使う書式（必要最小限の便利マクロ） --------------
% ベクトル表記：\vect{x} → 太字 x
\newcommand{\vect}[1]{\mathbf{#1}}

% （任意）よく使う集合：\R, \N
% \newcommand{\R}{\mathbb{R}}
% \newcommand{\N}{\mathbb{N}}

% ---------------- 執筆ルールTIPS（参照ラベル命名規約の提案） ----------
% ラベルは種類ごとに接頭辞を付けると検索しやすい：
%   図     : \label{fig:setup}
%   表     : \label{tab:hyperparam}
%   式     : \label{eq:snr}
%   節/章  : \label{sec:methods}, \label{chap:results}
% 参照例：
%   \figref{fig:setup}に計測システムを示す．\tabref{tab:hyperparam}に計測結果をまとめる．
% =====================================================================